\section{Réalisation centrale des trajectoires et holonomie de Sommerfeld
}
\label{vi:sec:kepler}

La particule définie dans la section~\ref{vi:sec:particula} possède
une histoire observable avant de recevoir une réalisation orbitale.
La question de cette section est de savoir comment une représentation de cette
histoire conserve flux, aire et phase lorsqu’elle s’exprime comme
mouvement central. La réponse s’organise en quatre opérations :
annulation des flux intérieurs dans une coupure du graphe,
normalisation de sa frontière dans une coordinatisation métrique, évolution sous
l’opérateur central déclaré et transport de phase sur ses cycles.


L’ordre maintient visible quelle donnée intervient dans chaque affirmation.
Le graphe et ses registres proviennent d’APP--TRIT--TPK ; la coordinatisation
métrique spécifie comment sa frontière est mesurée ; la réalisation dynamique
fixe le lagrangien qui agit dans cette coordinatisation ; la condition semi-classique
utilise la constante réduite de Planck déjà construite et l’holonomie
de la droite transportée. Les résultats de géométrie orbitale sont
démontrés avec ces données de départ explicites. La mécanique
centrale est reconnue dans la réalisation résultante, sans faire de ses
constantes conventionnelles des entrées du générateur.


\subsection{Conservation du flux et lecture de la frontière
}
\label{vi:subsec:gauss}

Soit $G=(V,E)$ une coupure finie du graphe de routes de la
réalisation considérée, avec un arc dans chaque orientation de ses
arêtes. Un flux orienté est une fonction
$F:\vec E\to\mathbb R$ qui satisfait
$F(y,x)=-F(x,y)$. Pour un sommet $x$, sa divergence est

\begin{equation}
 \operatorname{div}F(x)
 =\sum_{y:\{x,y\}\in E}F(x,y).
 \label{vi:eq:divergencia-grafo}
\end{equation}
Pour $S\subseteq V$, le flux sortant s’évalue sur les
arcs d’origine dans $S$ et d’extrémité dans $V\setminus S$.


\begin{theorem}[Bilan de flux dans une coupure finie
]
\label{vi:thm:gauss}
Pour toute partie $S\subseteq V$,

\begin{equation}
 \sum_{x\in S}\operatorname{div}F(x)
 =\sum_{\substack{x\in S,\ y\notin S\\\{x,y\}\in E}}F(x,y).
 \label{vi:eq:gauss}
\end{equation}
\end{theorem}
\begin{proof}
Une arête dont les deux extrémités sont dans $S$ contribue deux fois à la
somme des divergences, au moyen de $F(x,y)+F(y,x)=0$. Les arêtes
dont les deux extrémités sont hors de $S$ ne contribuent pas. Il reste
exactement les arêtes qui traversent la frontière, avec
l’orientation indiquée.

\end{proof}

Le théorème transmet à la frontière le bilan produit à
l’intérieur. Pour une représentation cubique à voisins immédiats,
considérons

\[
 C_R=\{x\in\mathbb Z^3:\|x\|_\infty\leq R\},
 \qquad R\in\mathbb Z_{\geq0}.
\]
Chacune des six faces contient $(2R+1)^2$ liens
sortants dans sa direction normale. Les liens de directions
différentes sont comptés séparément, même lorsqu’ils partent
d’un même sommet d’arête ou de coin. Par conséquent,

\begin{equation}
 |\partial C_R|=6(2R+1)^2.
 \label{vi:eq:frontera-cubica}
\end{equation}
Ce recensement correspond aux liens sortants, non au nombre de
sommets de la frontière.


Si $Q_F$ est le flux total de la source enfermée, sa moyenne
par lien est $Q_F/[6(2R+1)^2]$. Sous la condition d’isotropie
moyenne du lecteur, cette moyenne est l’intensité attribuée à la
coquille. Définissons son rayon effectif, en unités du pas de
la coordinatisation, par

\begin{equation}
 4\pi r_{\mathrm{eff}}^2=6(2R+1)^2.
 \label{vi:eq:radio-efectivo}
\end{equation}
On obtient alors

\begin{equation}
 E_R=\frac{Q_F}{4\pi r_{\mathrm{eff}}^2}.
 \label{vi:eq:inverso-area}
\end{equation}
La coordonnée $\pi$ de cette normalisation est la sortie régionale
HMT déjà exprimée. L’identité provient du recensement et de la
définition du rayon ; le passage d’une moyenne à une réalisation
centrale utilise en outre l’isotropie indiquée. La coordinatisation
dimensionnelle décide ensuite si $Q_F$ représente un flux électrique,
gravitationnel ou une autre observable et comment elle convertit le pas
réticulaire en longueur. Le bilan fini conserve ce
statut distinct de chaque opération.


\subsection{Dynamique centrale dans le feuillet euclidien
}
\label{vi:subsec:dinamica-central}

Soit $\mu>0$ la masse réduite de la réalisation matérielle et soit
$\mathcal K>0$ son couplage central, de dimension
produit d’énergie et de longueur. Ces coordonnées sont reçues après
les constructions internes de masse, d’action et de couplage ;
leur origine et leur droite dimensionnelle sont conservées lors de la
réalisation. Le théorème suivant a pour domaine toutes
les réalisations qui, dans leur feuillet euclidien, adoptent

\begin{equation}
 L(\mathbf r,\dot{\mathbf r})
 =\frac\mu2\|\dot{\mathbf r}\|^2+\frac{\mathcal K}{r},
 \qquad
 H(\mathbf r,\mathbf p)
 =\frac{\|\mathbf p\|^2}{2\mu}-\frac{\mathcal K}{r},
 \quad r=\|\mathbf r\|>0,
 \quad\mathbf p=\mu\dot{\mathbf r}.
 \label{vi:eq:kepler-lh}
\end{equation}
L’exclusion de $r=0$ délimite le domaine régulier de ces
coordonnées. La trajectoire est considérée sur un intervalle de
temps sans collision.


\begin{proposition}[Plan orbital et équation de la conique
]
\label{vi:prop:conica}
Une solution non radiale de l’équation d’Euler--Lagrange de
\eqref{vi:eq:kepler-lh} conserve
$\boldsymbol\ell=\mathbf r\times\dot{\mathbf r}$, avec
$\ell=\|\boldsymbol\ell\|>0$. Dans le plan orthogonal à
$\boldsymbol\ell$, on peut écrire

\begin{equation}
 r(\theta)=\frac{p}{1+\epsilon\cos(\theta-\theta_0)},
 \qquad p=\frac{\mu\ell^2}{\mathcal K},
 \qquad \epsilon\geq0.
 \label{vi:eq:conica}
\end{equation}
L’énergie négative correspond au régime elliptique
$0\leq\epsilon<1$.

\end{proposition}
\begin{proof}
La dérivée du potentiel donne

\begin{equation}
 \mu\ddot{\mathbf r}=-\frac{\mathcal K}{r^3}\mathbf r.
 \label{vi:eq:fuerza-central}
\end{equation}
D’où,
$\dot{\boldsymbol\ell}=\mathbf r\times\ddot{\mathbf r}=0$.
La position et la vitesse restent orthogonales au vecteur
constant non nul $\boldsymbol\ell$, ce qui fixe le plan.
Dans ses coordonnées polaires, $r^2\dot\theta=\ell$. En introduisant
$u=1/r$ dans l’équation radiale, on obtient l’équation de Binet

\[
 u''(\theta)+u(\theta)=\frac{\mathcal K}{\mu\ell^2}.
\]
Sa solution générale est
$u=\mathcal K[1+\epsilon\cos(\theta-\theta_0)]/(\mu\ell^2)$,
et son inversion prouve \eqref{vi:eq:conica} dans le domaine
où le dénominateur est positif. Substituer la solution dans
l’énergie donne
$H=\mathcal K(\epsilon^2-1)/(2p)$, qui est négative
exactement lorsque $\epsilon<1$.

\end{proof}

La lettre $\epsilon$ désigne ici l’excentricité orbitale et
non le nombre d’Euler. Le vecteur $\boldsymbol\ell$ est un
moment cinétique spécifique, non la constante de Planck.
Ces choix de notation permettent de comparer, sans les mélanger,
les observables géométriques et les constantes internes qui
interviennent dans leur réalisation.


\begin{proposition}[Conservation aréale et loi de la période
]
\label{vi:prop:ley-areal}
La vitesse aréolaire de l’orbite est constante :

\begin{equation}
 \frac{d\mathcal A_{\mathrm{orb}}}{dt}=\frac\ell2.
 \label{vi:eq:ley-areal}
\end{equation}
Pour une ellipse de demi-grand axe $a$, sa période satisfait

\begin{equation}
 T^2=\frac{4\pi^2\mu}{\mathcal K}a^3.
 \label{vi:eq:tercera-ley}
\end{equation}
\end{proposition}
\begin{proof}
Le secteur balayé pendant le temps $dt$ a pour aire
$r^2d\theta/2=\ell\,dt/2$. Si $b=a\sqrt{1-\epsilon^2}$
est le demi-petit axe, l’aire totale est $\pi ab$ et
$T=2\pi ab/\ell$. La relation
$p=a(1-\epsilon^2)=\mu\ell^2/\mathcal K$ implique
$\ell^2=\mathcal K a(1-\epsilon^2)/\mu$.
Substituer dans $T^2=4\pi^2a^2b^2/\ell^2$ donne l’expression
\eqref{vi:eq:tercera-ley}.

\end{proof}

La conservation aréale est déterminée par la symétrie centrale
du mouvement. L’ellipse de cette proposition vit dans
l’espace de configuration. L’ellipse d’action construite
à partir de la paire angulaire vit, en revanche, dans sa coordinatisation d’action
et transporte les normalisations correspondantes. Le lien
entre les deux exige l’application de réalisation qui identifie leurs
variables ; partager une forme elliptique ne remplace pas cette application.
La lecture orbitale conserve ainsi les dépendances de masse,
de couplage et de géométrie qui la produisent.


\subsection{Revêtement orbital et droite réelle de signe
}
\label{vi:subsec:cubierta-orbital}

Dans le plan complexe privé de l’origine, l’application

\begin{equation}
 q=z^2:\mathbb C\setminus\{0\}\longrightarrow
                 \mathbb C\setminus\{0\}
 \label{vi:eq:levi-civita}
\end{equation}
est un revêtement double. Si $q(t)=r(t)e^{i\theta(t)}$ et si l’on
choisit initialement une racine, son relèvement est
$z(t)=\sqrt{r(t)}e^{i\theta(t)/2}$. Un tour avec un incrément
$2\pi$ de $\theta$ se termine en $-z(0)$, et deux tours
restituent le point initial. C’est le transport des deux
relèvements sur le plan perforé.


Une bande orbitale requiert un autre objet : une droite réelle
associée à l’orbite et à son caractère de signe. Soit
$\gamma:S^1\to\mathbb R^2$ un paramétrage régulier
injectif d’une orbite elliptique et soit $\delta>0$. Avec
l’holonomie $-1$ en un tour, sa bande d’intervalles se
représente par

\begin{equation}
 M_\gamma=([0,2\pi]\times[-\delta,\delta])/
          \bigl((0,u)\sim(2\pi,-u)\bigr).
 \label{vi:eq:mobius-orbital}
\end{equation}
L’identification de la base à l’image de $\gamma$
situe la section nulle sur l’orbite.


\begin{proposition}[Transport orienté sur le cycle orbital
]
\label{vi:prop:holonomia-orbital}
Le transport de la droite de \eqref{vi:eq:mobius-orbital}
satisfait
$\mathcal H(\gamma)=-I$ et $\mathcal H(\gamma^2)=I$.

\end{proposition}
\begin{proof}
La fibre terminale s’identifie à l’initiale par
$u\mapsto-u$ ; composer cette application deux fois donne l’identité.
C’est la construction du
théorème~\ref{vi:thm:retorno-mobius}, appliquée à la boucle orbitale.

\end{proof}

Le revêtement \eqref{vi:eq:levi-civita} possède des fibres de deux
points ; \eqref{vi:eq:mobius-orbital} possède des fibres
d’intervalles. La monodromie d’ordre deux apparaît dans les deux,
mais conserve ces types distincts. La régularisation
dynamique de Levi--Civita utilise en outre la transformation
des moments et le reparamétrage temporel. Dans la présente
construction, on utilise uniquement son revêtement et sa loi de
relèvement, avec la portée qui vient d’être démontrée.


\subsection{Condition de phase et déplacement demi-entier
}
\label{vi:subsec:sommerfeld}

Soit $\gamma$ un cycle de la réalisation hamiltonienne et soit

\begin{equation}
 J_\gamma=\oint_\gamma\mathbf p\cdot d\mathbf q
 \label{vi:eq:accion-ciclo}
\end{equation}
son action. La constante réduite de Planck
$\hbar_{\mathrm{int}}>0$ est reçue de la construction HMT
d’action, dans la même droite dimensionnelle que $J_\gamma$.
Le quotient $J_\gamma/\hbar_{\mathrm{int}}$ est donc
adimensionnel. Supposons que le transport semi-classique dans la
droite complexifiée apporte la phase
$\exp(iJ_\gamma/\hbar_{\mathrm{int}})$ et que la droite réelle
d’orientation apporte le caractère
$\varepsilon_\gamma\in\{+1,-1\}$. La condition de fermeture
de la section transportée est

\begin{equation}
 \varepsilon_\gamma
       \exp\!\left(\frac{iJ_\gamma}{\hbar_{\mathrm{int}}}\right)=1.
 \label{vi:eq:cierre-sommerfeld}
\end{equation}
Le produit conserve les deux contributions : action du
cycle et signe de la droite. La seconde n’est pas absorbée dans une
modification de la constante de Planck.


\begin{theorem}[Secteurs entier et demi-entier de l’action
]
\label{vi:thm:sommerfeld}
Sous le transport \eqref{vi:eq:cierre-sommerfeld},

\begin{equation}
 J_\gamma=
 \begin{cases}
  2\pi\hbar_{\mathrm{int}}n,
          &\varepsilon_\gamma=+1,\\[1mm]
  2\pi\hbar_{\mathrm{int}}(n+\tfrac12),
          &\varepsilon_\gamma=-1,
 \end{cases}
 \qquad n\in\mathbb Z.
 \label{vi:eq:accion-sommerfeld}
\end{equation}
\end{theorem}
\begin{proof}
Pour un signe positif, l’exponentielle doit valoir $1$, dont les
arguments réels sont $2\pi n$. Pour un signe négatif, elle doit
valoir $-1$, dont les arguments sont $(2n+1)\pi$. Multiplier
par $\hbar_{\mathrm{int}}$ produit les deux collections.

\end{proof}

La formule détermine la restriction de fermeture sur
l’action d’un cycle donné. Pour obtenir un spectre atomique,
on utilise en outre l’hamiltonien spécifique, les cycles
indépendants et la prescription de quantification. Si cette
prescription comprend des indices de Maslov, leurs phases se
composent avec \eqref{vi:eq:cierre-sommerfeld} conformément au
transport choisi, en évitant de compter deux fois une même
contribution topologique.


Inverser l’orientation du cycle remplace $J_\gamma$ par
$-J_\gamma$, tandis que le caractère de signe satisfait
$\varepsilon_{\gamma^{-1}}=\varepsilon_\gamma^{-1}
=\varepsilon_\gamma$. La phase totale se transforme en son
inverse et la condition de fermeture est conservée. C’est une
loi précise de bidirectionnalité orbitale. Son identification
aux antisections matérielles requiert le relèvement
de la section~\ref{vi:subsec:antisecciones} ; l’égalité
de deux signes n’établit pas à elle seule cette identification.


La réalisation réunit ainsi trois conservations
consécutives : le flux passe de l’intérieur à la frontière ;
la dynamique centrale conserve le moment cinétique et l’aire
balayée par unité de temps ; le transport de la section
conserve sa condition de fermeture au moyen de l’action et de
l’holonomie. Chaque résultat conserve l’équation, le domaine
et les données qui le soutiennent. L’histoire matérielle du
TPK est représentée dans cette dynamique avec sa mémoire, au lieu
d’être remplacée par la figure orbitale qu’elle exprime.

